\documentclass[a4paper, 11pt]{article}

% Packages
\usepackage[utf8]{inputenc}
\usepackage[T1]{fontenc}
\usepackage{lmodern}
\usepackage[french]{babel}
\usepackage{textcomp} % Required for some symbols
\usepackage[utf8]{inputenc} % Specify input encoding
\usepackage{amsmath}
\usepackage{amssymb}
\usepackage{geometry}
\usepackage{hyperref}
\usepackage{algorithm}
\usepackage[noend]{algpseudocode}
\usepackage{listings}
\usepackage{xcolor}

% Listings settings
\definecolor{codegreen}{rgb}{0,0.6,0}
\definecolor{codegray}{rgb}{0.5,0.5,0.5}
\definecolor{codepurple}{rgb}{0.58,0,0.82}
\definecolor{backcolour}{rgb}{0.95,0.95,0.92}

\lstdefinestyle{mystyle}{
    inputencoding=utf8, % Handle UTF-8 characters
    backgroundcolor=\color{backcolour},   
    commentstyle=\color{codegreen},
    keywordstyle=\color{magenta},
    numberstyle=\tiny\color{codegray},
    stringstyle=\color{codepurple},
    basicstyle=\footnotesize\ttfamily,
    breakatwhitespace=false,         
    breaklines=true,                 
    captionpos=b,                    
    keepspaces=true,                 
    numbers=left,                    
    numbersep=5pt,                  
    showspaces=false,                
    showstringspaces=false,
    showtabs=false,                  
    tabsize=2,
    language=C
}
\lstset{style=mystyle}

% Document settings
\geometry{a4paper, top=2.5cm, bottom=2.5cm, left=2.5cm, right=2.5cm}
\hypersetup{
    colorlinks=true,
    linkcolor=blue,
    filecolor=magenta,      
    urlcolor=cyan,
}

\title{\textbf{Compte Rendu : Structures de Données Avancées}}
\author{Module : Algorithmique et Programmation \\ Année : 1ère Année ILISI}
\date{\today}

\begin{document}

\maketitle
\hrule

\section*{Introduction}

Ce document présente l'analyse, la conception et l'implémentation de plusieurs structures de données, conformément aux exigences du travail pratique. Il est divisé en deux problèmes principaux :
\begin{enumerate}
    \item L'implémentation d'une \textbf{file} et d'une \textbf{liste} à l'aide d'un tableau circulaire.
    \item L'implémentation de \textbf{deux piles} dans un seul et même tableau.
\end{enumerate}
Pour chaque problème, nous suivrons une démarche structurée : analyse du problème, proposition d'une solution, description des algorithmes et enfin, présentation du programme final avec une démonstration de son fonctionnement.

\hrule\vspace{1cm}

\section{Problème 1 : Le Tableau Circulaire}

\subsection{I. Analyse du Problème}
Le premier problème nous demande d'utiliser le concept de \textbf{tableau circulaire} pour implémenter une \textbf{File} (structure de type FIFO). L'objectif est de gérer un espace mémoire fixe de manière efficace.

\subsection{II. Solution Proposée}
Pour la file, nous utilisons un tableau \texttt{buffer} et trois variables de contrôle : \texttt{tete}, \texttt{queue}, et \texttt{taille}. La structure est définie comme suit :
\begin{lstlisting}
typedef struct FileCirculaire
{
    int buffer[MAX];
    int queue; // La queue, d'ou on defile
    int tete;  // La tete, ou on enfile
    int taille;
} FileCirculaire;
\end{lstlisting}

\subsection{III. Algorithmes et Exemples}
Une file est une structure de données de type \textbf{premier entré, premier sorti (FIFO)}.

\subsubsection*{\texttt{enfiler(file, valeur)}}
L'ajout se fait à la \texttt{tete}.
\begin{algorithmic}
\If{file.taille = MAX}
    \State Afficher "Erreur: File pleine"
\Else
    \State file.buffer[file.tete] $\leftarrow$ valeur
    \State file.tete $\leftarrow$ (file.tete + 1) \textbf{mod} MAX
    \State file.taille $\leftarrow$ file.taille + 1
\EndIf
\end{algorithmic}
Exemple d'utilisation :
\begin{lstlisting}
FileCirculaire *file = initialiser_file();
enfiler(file, 10); // Ajoute 10
enfiler(file, 20); // Ajoute 20
\end{lstlisting}

\subsubsection*{\texttt{defiler(file)}}
Le retrait se fait à la \texttt{queue}.
\begin{algorithmic}
\If{file.taille = 0}
    \State Afficher "Erreur: File vide"
\Else
    \State element $\leftarrow$ file.buffer[file.queue]
    \State file.queue $\leftarrow$ (file.queue + 1) \textbf{mod} MAX
    \State file.taille $\leftarrow$ file.taille - 1
    \State \Return element
\EndIf
\end{algorithmic}

\subsection{IV. Programme et Démonstration}
Le code source complet se trouve en annexe. La fonction \texttt{demonstration\_file\_circulaire()} montre son fonctionnement.

\section{Problème 2 : Deux Piles dans un Tableau}

\subsection{I. Analyse du Problème}
Le second problème consiste à implémenter deux piles distinctes en n'utilisant qu'un seul tableau pour optimiser l'utilisation de l'espace.

\subsection{II. Solution Proposée}
La solution consiste à faire croître les deux piles en sens opposé. La pile gauche part de l'indice \texttt{0} et la pile droite de \texttt{MAX-1}. La structure est :
\begin{lstlisting}
typedef struct DoublePile
{
    int tab[MAX];
    int sommet_g; // Sommet de la pile gauche
    int sommet_d; // Sommet de la pile droite
} DoublePile;
\end{lstlisting}

\subsection{III. Algorithmes et Exemples}
Une pile est une structure de données de type \textbf{dernier entré, premier sorti (LIFO)}.

\subsubsection*{\texttt{empiler\_g(pile, valeur)}}
\begin{algorithmic}
\If{pile.sommet\_g + 1 = pile.sommet\_d}
    \State Afficher "Erreur: Piles pleines"
\Else
    \State pile.sommet\_g $\leftarrow$ pile.sommet\_g + 1
    \State pile.tab[pile.sommet\_g] $\leftarrow$ valeur
\EndIf
\end{algorithmic}
Exemple d'utilisation :
\begin{lstlisting}
DoublePile *pile = initialiser_dp();
empiler_g(pile, 5);  // Ajoute 5 a gauche
empiler_d(pile, 99); // Ajoute 99 a droite
\end{lstlisting}

\subsection{IV. Programme et Démonstration}
Le code de la \texttt{DoublePile} est également en annexe. La fonction \texttt{demonstration\_double\_pile()} illustre son utilisation.

\newpage
\appendix
\section{Code Source Complet : \texttt{main.c}}
\lstinputlisting{main.c}

\end{document}
